\begin{figure}[H]
\centering
\includegraphics[width=1.0\linewidth]{wakayama/F1_alt_half_speed.pdf}
\caption[時期×高度の風速と風向]{上：平均風速。下：平均風向（吹いてくる向き）。風向には円環用のカラーマップを使っている。順次カラーマップを使うと \ang{0} と \ang{360} の境目に実在しない断層が現れるためである。灰色は定常度が 0.3 未満で塗りを外した領域で、そのビンではそもそも風向が定まっていない。}
\label{fig:f1-wakayama}
\end{figure}
